% preamble.tex -- the preamble of reports/status-2026-11, with: a glyph fallback for the symbols
% the programme's files carry (STIX Two Math, then DejaVu Sans), \tid{} for target ids, and the
% PDF's optional info suppressed so that two builds give the same bytes (see the Makefile).
\documentclass[11pt,a4paper]{article}

\usepackage[margin=2.2cm]{geometry}
\usepackage{amsmath,amssymb,amsthm}
\usepackage{fontspec}
\usepackage{unicode-math}
\directlua{luaotfload.add_fallback("progfallback",
  {"STIXTwoMath-Regular.otf:mode=node;", "DejaVuSans.ttf:mode=node;"})}
\setmainfont{STIX Two Text}[RawFeature={fallback=progfallback}]
\setmathfont{STIX Two Math}
\setmonofont{DejaVu Sans Mono}[Scale=0.80, RawFeature={fallback=progfallback}]
\usepackage{booktabs,longtable,array}
\usepackage{fancyvrb}
\usepackage{luacode}
\usepackage{enumitem}
\usepackage[hidelinks]{hyperref}
\setlist{itemsep=1pt,topsep=3pt}
\setlength{\emergencystretch}{3em}

% reproducible output: no PTEX.FullBanner, no PTEX.FileName (the dates and the /ID follow
% SOURCE_DATE_EPOCH, which the Makefile sets to the commit time of REV)
\pdfvariable suppressoptionalinfo \numexpr 1 + 2 \relax

% ---------------------------------------------------------------- Lean names
% \lean{name}: typeset a Lean identifier verbatim in monospace, with break points after '_', '.',
% ':' and '/', so long names do not overflow table cells. \tid{T-...}: a target id, breakable
% after '-' as well.
\begin{luacode*}
function leanname(s, breaks)
  breaks = breaks or "_.:/"
  local buf = ""
  local prev = ""
  for _, c in utf8.codes(s) do
    local ch = utf8.char(c)
    if prev:match("^%l$") and ch:match("^%u$") then
      tex.sprint(-2, buf); buf = ""
      tex.sprint("\\penalty100\\hskip0pt\\relax{}")
    end
    buf = buf .. ch
    if breaks:find(ch, 1, true) then
      tex.sprint(-2, buf); buf = ""
      tex.sprint("\\allowbreak{}")
    end
    prev = ch
  end
  tex.sprint(-2, buf)
end
\end{luacode*}
\newcommand{\lean}[1]{\texttt{\directlua{leanname("\luaescapestring{\detokenize{#1}}")}}}
\newcommand{\tid}[1]{\texttt{\directlua{leanname("\luaescapestring{\detokenize{#1}}", "-_.:/")}}}
\newcommand{\lbl}[1]{\lean{#1}}
\newcommand{\file}[1]{\lean{#1}}
